\chapter{Motiva\c tie}
\sloppy


\section{Contextul \c si problema}

Organizarea \textbf{conferin\c telor, meeting-urilor, concertelor} \c si a altor evenimente cu public \c si locuri limitate implic\u a, \^in mod tradi\c tional, procese fragmentate: promovarea prin canale diferite (re\c tele sociale, email, site-uri proprii), v\^anzarea biletelor sau \^inscrierilor \^in format fizic sau prin transfer bancar manual, \c si validarea participan\c tilor la intrare prin liste pe h\^artie. Aceast\u a abordare genereaz\u a erori, pierderi de timp \c si dificult\u a\c ti \^in reconcilierea pl\u a\c tilor cu locurile rezervate.

Digitalizarea procesului — centralizarea informa\c tiilor despre evenimente, automatizarea pl\u a\c tilor \c si asocierea fiec\u arui bilet cu identitatea participantului — reprezint\u a un pas natural \^in contextul extinderii serviciilor online \c si al a\c stept\u arilor utilizatorilor privind comoditatea \c si securitatea tranzac\c tiilor \cite{fielding2000, stripe2024}.

Organizatorii (universit\u a\c ti, firme, asocia\c tii culturale, s\u ali de concert) nu dispun adesea de resurse pentru solu\c tii enterprise costisitoare. \^In acela\c si timp, platformele interna\c tionale generice nu sunt \^intotdeauna adaptate pie\c tei locale: moneda RON, structura administrativ\u a pe jude\c te, sau necesitatea unui flux simplu de check-in administrat de organizator. Pentru \textbf{evenimente sportive} amatoriale, acelea\c si mecanisme (\^inscriere, plat\u a, list\u a nominal\u a) sunt utile, dar reprezint\u a un caz secundar fa\c t\u a de conferin\c te \c si concerte.

Participan\c tii au nevoie de o interfa\c t\u a clar\u a: descoperire evenimente, filtrare, achizi\c tia rapid\u a a unuia sau  a mai multor bilete nominale, vizualizarea istoricului \c si gestionarea contului. Organizatorii au nevoie de un panou unic pentru publicare, monitorizare locuri disponibile \c si confirmare prezen\c t\u a la fa\c ta locului.




\section{Obiectivele lucr\u arii}

Obiectivul principal este proiectarea \c si implementarea \textbf{Aplica\c tiei web pentru organizarea \c si gestionarea evenimentelor}, care s\u a \^indeplineasc\u a urm\u atoarele obiective specifice:

\begin{enumerate}
	\item Implementarea unui modul de autentificare securizat: \^inregistrarea cu verificarea e-mailului, login JWT, resetarea parolei, editarea profilului \c si \c stergerea contului.
	\item Oferirea unei interfe\c te publice pentru listarea evenimentelor viitoare, cu filtre dup\u a nume, jude\c t \c si dat\u a.
	\item Permiterea achizi\c tion\u arii biletelor nominale (unul sau mai multe nume per tranzac\c tie) cu plat\u a online prin Stripe Checkout, \^in moneda RON.
	\item Dezvoltarea unui panou de administrare pentru CRUD evenimente (inclusiv imagini), list\u a participan\c ti \c si check-in.
	\item Gestionarea rolurilor utilizator/admin \c si promovarea/revocarea administratorilor cu parol\u a de autorizare.
	\item Documentarea arhitecturii, a modelului de date \c si a fluxurilor principale (UML, ERD).
\end{enumerate}

Lucrarea demonstreaz\u a integrarea tehnologiilor web moderne \^intr-un produs coerent, testabil \c si prezentabil, cu respectarea unor practici de securitate recunoscute \cite{fowler2002}.


\section{Structura documentului}

Prezenta lucrare este organizat\u a structural pe capitole, dup\u a cum urmeaz\u a:

\vspace{0.2cm}
\begin{description}
	\setlength{\itemsep}{0.3cm}
	\item[Capitolul 2] Analizeaz\u a \^in detaliu starea actual\u a a domeniului. Cuprinde studii relevante din literatur\u a, o analiz\u a a platformelor existente pe pia\c t\u a \c si o compara\c tie direct\u a cu solu\c tia propus\u a.
	\item[Capitolul 3] Descrie \^in profunzime arhitectura sistemului \c si detaliile de implementare ale aplica\c tiei. Sunt incluse cerin\c tele func\c tionale \c si non-func\c tionale, tehnologiile alese, diagramele UML, structura bazei de date \c si fragmente de cod surs\u a relevante.
	\item[Capitolul 4] Prezint\u a interfe\c tele grafice ale aplica\c tiei prin intermediul unor capturi de ecran detaliate, mapate pe fluxurile principale de utilizare.
	\item[Capitolul 5] Sintetizeaz\u a concluziile generale ale proiectului, eviden\c tiaz\u a limit\u arile identificate pe parcursul dezvolt\u arii \c si propune viitoarele direc\c tii de dezvoltare.
\end{description}


\section{Tipuri de evenimente vizate}
\sloppy
Aplica\c tia web pentru organizarea \c si gestionarea evenimentelor este conceput\u a \^in jurul unor formate frecvente \^in via\c ta academic\u a, cultural\u a \c si profesional\u a. \textit{Tabelul 1.1} sintetizeaz\u a ca-tegoriile principale, actorii implica\c ti \c si modul \^in care platforma le sus\c tine. 

\begin{table}[htbp]
	\centering
	\renewcommand{\arraystretch}{1.5} 
	\setlength{\tabcolsep}{10pt} %	
	\begin{tabular}{|l|l|p{5.5cm}|} 
		\hline
		\textbf{Tip Eveniment} & \textbf{Actori Principali} & \textbf{Scenariu \^in Aplica\c tie} \\ \hline
		Conferin\c t\u a / Meeting & Organizator, Participant & Publicare cu detalii complete (dat\u a, locuri limitate, descriere, pre\c t \^in RON), \^inregistrare online \c si check-in la intrare. \\ \hline
		Concert / Festival & Organizator, Public & Accent pe componenta vizual\u a (galerie de imagini), redirec\c tionare rapid\u a c\u atre Stripe Checkout pentru achizi\c tie bilet nominal. \\ \hline
		Competi\c tie Sportiv\u a & Admin, Sportiv & \^Inscriere pe baz\u a de bilet nominal, check-in realizat de organizator la \^inregistrarea \^in concurs (caz secundar). \\ \hline
	\end{tabular}

	\caption{Tipuri de evenimente reflectate \^in aplica\c tie}
	\label{tab:evenimente}
\end{table}

\subsection{Conferin\c te, meeting-uri \c si evenimente profesionale}

Universit\u a\c ti, camere de comer\c t \c si firme de consultan\c t\u a organizeaz\u a periodic conferin\c te \c si meeting-uri cu num\u ar limitat de locuri. Participan\c tii trebuie s\u a se \^inregistreze, s\u a-\c si confirme identitatea (email verificat) \c si, uneori, s\u a achite o tax\u a de participare. Platforma acoper\u a acest flux end-to-end: publicarea evenimentului cu dat\u a, loca\c tie (jude\c t, localitate) \c si capacitate, urmat\u a de \^inregistrare online \c si emitere bilet nominal.

Pentru meeting-uri interne sau seminarii pl\u atite, administratorul poate crea evenimente cu pre\c t fix \c si num\u ar redus de locuri; sistemul decrementeaz\u a automat locurile disponibile la fiecare tranzac\c tie confirmat\u a. Check-in-ul din panoul admin permite validarea prezen\c tei la recep\c tie sau la intrarea \^in sal\u a util at\^at pentru conferin\c te, c\^at \c si pentru training-uri corporate.

\subsection{Concerte \c si evenimente culturale}

Concertele \c si festivalurile locale necesit\u a o prezentare atractiv\u a (titlu, descriere, mai multe imagini) \c si v\^anzare rapid\u a a biletelor. Aplica\c tia permite \^inc\u arcarea a p\^an\u a la opt imagini per eveniment, afi\c sarea galeriei pe pagina de detalii \c si redirect c\u atre Stripe Checkout pentru plat\u a securizat\u a.

\newpage Filtrarea dup\u a jude\c t \c si dat\u a ajut\u a publicul s\u a descopere evenimente din apropiere relevant pentru turnee na\c tionale de concerte sau festivaluri itinerante. Biletele nominale asigur\u a c\u a fiecare loc este asociat unui participant, reduc\^and rev\^anzarea neautorizat\u a \c si simplific\^and controlul la intrare.

\subsection{Evenimente sportive \c si utilizare complementar\u a}

Modelul actual al aplica\c tiei (nume pe bilet, check-in administrat manual) se preteaz\u a \c si competi\c tiilor sportive amatoriale sau turneelor locale, unde organizatorul trebuie s\u a verifice identitatea participantului. Aceast\u a ni\c s\u a este secundar\u a fa\c t\u a de conferin\c te, concerte \c si meeting-uri: nu sunt implementate func\c tionalit\u a\c ti specifice sportului (categorii de v\^arst\u a, echipe, program meciuri), dar fluxul general de \^inscriere \c si plat\u a r\u am\^ane valabil.

Extinderi viitoare (categorii de bilete, abonamente de sezon) ar putea \^int\u ari suportul sportiv f\u ar\u a a modifica arhitectura de baz\u a.

\subsection{ Beneficii transversale}



Indiferent de tipul evenimentului, organizatorul beneficiaz\u a de: 

\begin{itemize}
	\item \textbf{Un singur panou} pentru toate formatele: nu este nevoie de solu\c tii separate pentru conferin\c te vs. concerte.
	\item \textbf{Trasabilitate pl\u a\c ti}: fiecare bilet este legat de o sesiune Stripe \c si de un utilizator a-utentificat.
	\item \textbf{Adaptare local\u a}: jude\c te rom\^ane\c sti, moned\u a RON, e-mailuri prin SMTP Brevo.
	\item \textbf{Control asupra datelor}: baza PostgreSQL self-hosted, f\u ar\u a dependen\c t\u a de un marketplace global.
\end{itemize}

\vspace{1em} % Spa\c tiu \^intre list\u a \c si urm\u atorul paragraf

Aceast\u a diversitate de cazuri de utilizare justific\u a alegerea unui model de date simplu (eveniment unic, pre\c t unic, bilete nominale) care r\u am\^ane suficient de expresiv pentru majoritatea scenariilor \^int\^alnite.